% DEFINED:   sec:book-close
% RESOLVED:  sec:three-axes, sec:inductive-biases, ch:foundations
%            (all backward). NO \Cref to Chapter 7 (the journey narrates the
%            causal trace in prose, but ties back to the supervised arc).
% FORWARD:   NONE. This is the last section of the book.

\section{Closing: the book, completed}\label{sec:book-close}

This is where the book ends, so it is worth saying plainly what the eight
chapters were for. The thread was a single equation with three slots: a
parametric model, an output distribution, and a loss, trained by gradient
descent on data. Choosing what to put in those slots is choosing an inductive
bias, and the whole book was a tour of those choices along three axes
(\Cref{sec:three-axes}).

\begin{enumerate}
\item \textbf{Architecture.} Structural priors about how features compose:
locality and weight-sharing in a convolutional network, recurrence in a
recurrent one, a bounded window in the fixed-context model, content-addressable
lookup in a transformer. Each is a statement about which functions are easy to
represent and which the data should not have to teach from scratch.
\item \textbf{Output head.} Priors about the distribution of the response:
Bernoulli for a bit, categorical for a class, Gaussian for a real value,
Poisson for a count. The head fixes what ``being right'' even means, and the
loss follows from it.
\item \textbf{Implementation realization.} Whether gradient descent, from a
given initialization and schedule, actually finds one of the bias-permitted
solutions, and whether we can read that solution back off the trained weights.
\end{enumerate}

The journey ran the length of all three. It opened (\Cref{ch:foundations}) with
the smallest non-trivial case, a multilayer perceptron learning XOR, the
problem a single linear boundary provably cannot solve and a hidden layer can:
the first inductive bias in the book (\Cref{sec:inductive-biases}), the bias
toward composing simple boundaries into a curved one. It then walked the output
head through its distributional choices, and the architecture through the
convolutional, recurrent, fixed-context, and attention families, each matched
to a structure in its data. The pointer transformer was the high-water mark of
the architecture axis, a model that solves content-addressable lookup because
the lookup is built into the attention operation. Then the book turned the lens
around and read the bias back out of a trained model, finding that having the
right architecture is necessary but not sufficient: at the smallest scale the
learned circuit was perfectly legible, and one increment of scale later a
causal trace, not the attention maps, was what recovered the mechanism. That
was the third axis made empirical, and the book's most honest result, that what
gradient descent realized is a separate question from what the architecture
permits.

This final chapter changed the learning signal one last time. The per-example
label disappeared and came back as a scalar reward over a whole trajectory, so
the hard new problem was temporal credit assignment. And yet the inner loop
turned out to be the toolkit from the very first chapter: REINFORCE is softmax
cross-entropy on the actions the agent took, the predicted distribution minus
the one-hot chosen action, scaled by the return. The supervised gradient of
\Cref{ch:foundations} is, line for line, what a policy-gradient agent runs;
what reinforcement learning adds is the outer loop, the credit-assignment
machinery, and the trajectory distribution over which that inner loss is
weighted. The same is true above the gridworld: softmax cross-entropy is the
gradient signal in behavioral cloning, in policy gradient, in the clipped
surrogate of proximal policy optimization, in the policy network of AlphaZero;
squared error is the gradient signal in value-function regression and in the
reward-model fitting step of learning from human feedback. The supervised
toolkit is the inner loop of every modern learning system; the rest is which
labels you feed it and how you weight them.

Two incomputable ideals bracketed the book. Solomonoff induction was the
optimal predictor that no machine can run, and the language models of Part~II
were its computable approximations. AIXI is its action-taking sibling, the
optimal agent that no machine can run, and the REINFORCE agent of this chapter
is its smallest computable shadow: a single policy network, a single return, a
single gradient step. Between the incomputable optimum and the runnable
approximation sits the architecture, which is exactly where the compromise
lives, and which is exactly what an inductive bias is. The boundary between
supervised learning and reinforcement learning, like the boundary between
prediction and action, turns out to be conceptual rather than categorical, and
the systems that matter most today thread back and forth across it. The math
that carries across, the math worked out by hand across these eight chapters,
is the part worth keeping. That is the end of the book.
